\begin{thebibliography}{99}

\bibitem{apa2013} American Psychiatric Association, \emph{Diagnostic and Statistical Manual of
Mental Disorders}, 5th ed. Arlington, VA, USA: American Psychiatric Publishing, 2013.

\bibitem{polanczyk2007} G. Polanczyk, M. S. de Lima, B. L. Horta, J. Biederman, and L. A.
Rohde, ``The worldwide prevalence of ADHD: A systematic review and metaregression analysis,''
\emph{Am. J. Psychiatry}, vol.~164, no.~6, pp.~942--948, 2007.

\bibitem{maulik2011} P. K. Maulik, M. N. Mascarenhas, C. D. Mathers, T. Dua, and S. Saxena,
``Prevalence of intellectual disability: A meta-analysis of population-based studies,''
\emph{Res. Dev. Disabil.}, vol.~32, no.~2, pp.~419--436, 2011.

\bibitem{chandela2025} U. Chandela, K. N. Faisal, and R. R. Sharma, ``Electroencephalogram-based
unified approach for multiple neurodevelopmental disorders detection in children using
successive multivariate variational mode decomposition,'' \emph{IEEE Trans. Cogn. Develop.
Syst.}, vol.~17, no.~6, pp.~1350--1359, Dec. 2025.

\bibitem{sareen2020data} E. Sareen, L. Singh, B. Varkey, K. Achary, and A. Gupta, ``EEG
dataset of individuals with intellectual and developmental disorder and healthy controls
under rest and music stimuli,'' \emph{Data Brief}, vol.~30, 2020, Art. no.~105488.

\bibitem{motie2020} A. Motie Nasrabadi, A. Allahverdy, M. Samavati, and M. R. Mohammadi,
``EEG data for ADHD/control children,'' IEEE DataPort, 2020, doi: 10.21227/rzfh-zn36.

\bibitem{snyder2006} S. M. Snyder and J. R. Hall, ``A meta-analysis of quantitative EEG power
associated with attention-deficit hyperactivity disorder,'' \emph{J. Clin. Neurophysiol.},
vol.~23, no.~5, pp.~440--455, 2006.

\bibitem{arns2013} M. Arns, C. K. Conners, and H. C. Kraemer, ``A decade of EEG theta/beta
ratio research in ADHD: A meta-analysis,'' \emph{J. Atten. Disord.}, vol.~17, no.~5,
pp.~374--383, 2013.

\bibitem{clarke2002} A. R. Clarke, R. J. Barry, D. Bond, R. McCarthy, and M. Selikowitz,
``Effects of stimulant medications on the EEG of children with attention-deficit/hyperactivity
disorder,'' \emph{Psychopharmacology}, vol.~164, no.~3, pp.~277--284, 2002.

\bibitem{huang1998} N. E. Huang \emph{et al.}, ``The empirical mode decomposition and the
Hilbert spectrum for nonlinear and non-stationary time series analysis,'' \emph{Proc. R. Soc.
Lond. A}, vol.~454, pp.~903--995, 1998.

\bibitem{dragomiretskiy2014} K. Dragomiretskiy and D. Zosso, ``Variational mode
decomposition,'' \emph{IEEE Trans. Signal Process.}, vol.~62, no.~3, pp.~531--544, Feb. 2014.

\bibitem{rehman2019} N. Ur Rehman and H. Aftab, ``Multivariate variational mode
decomposition,'' \emph{IEEE Trans. Signal Process.}, vol.~67, no.~23, pp.~6039--6052, Dec. 2019.

\bibitem{liu2022} S. Liu and K. Yu, ``Successive multivariate variational mode
decomposition,'' \emph{Multidimensional Syst. Signal Process.}, vol.~33, no.~3,
pp.~917--943, 2022.

\bibitem{peng2005} H. Peng, F. Long, and C. Ding, ``Feature selection based on mutual
information: Criteria of max-dependency, max-relevance, and min-redundancy,'' \emph{IEEE
Trans. Pattern Anal. Mach. Intell.}, vol.~27, no.~8, pp.~1226--1238, Aug. 2005.

\bibitem{cover1967} T. Cover and P. Hart, ``Nearest neighbor pattern classification,''
\emph{IEEE Trans. Inf. Theory}, vol.~13, no.~1, pp.~21--27, 1967.

\bibitem{welch1967} P. Welch, ``The use of fast Fourier transform for the estimation of power
spectra: A method based on time averaging over short, modified periodograms,'' \emph{IEEE
Trans. Audio Electroacoust.}, vol.~15, no.~2, pp.~70--73, 1967.

\bibitem{hjorth1970} B. Hjorth, ``EEG analysis based on time domain properties,''
\emph{Electroencephalogr. Clin. Neurophysiol.}, vol.~29, no.~3, pp.~306--310, 1970.

\bibitem{richman2000} J. S. Richman and J. R. Moorman, ``Physiological time-series analysis
using approximate entropy and sample entropy,'' \emph{Am. J. Physiol. Heart Circ. Physiol.},
vol.~278, no.~6, pp.~H2039--H2049, 2000.

\bibitem{bandt2002} C. Bandt and B. Pompe, ``Permutation entropy: A natural complexity measure
for time series,'' \emph{Phys. Rev. Lett.}, vol.~88, no.~17, Art. no.~174102, 2002.

\bibitem{donoghue2020} T. Donoghue \emph{et al.}, ``Parameterizing neural power spectra into
periodic and aperiodic components,'' \emph{Nat. Neurosci.}, vol.~23, pp.~1655--1665, 2020.

\bibitem{vinck2011} M. Vinck, R. Oostenveld, M. van Wingerden, F. Battaglia, and C. M. A.
Pennartz, ``An improved index of phase-synchronization for electrophysiological data in the
presence of volume-conduction, noise and sample-size bias,'' \emph{NeuroImage}, vol.~55, no.~4,
pp.~1548--1565, 2011.

\bibitem{chen2010} Y. Chen, A. Wiesel, Y. C. Eldar, and A. O. Hero, ``Shrinkage algorithms for
MMSE covariance estimation,'' \emph{IEEE Trans. Signal Process.}, vol.~58, no.~10,
pp.~5016--5029, 2010.

\bibitem{barachant2012} A. Barachant, S. Bonnet, M. Congedo, and C. Jutten, ``Multiclass
brain--computer interface classification by Riemannian geometry,'' \emph{IEEE Trans. Biomed.
Eng.}, vol.~59, no.~4, pp.~920--928, 2012.

\bibitem{congedo2017} M. Congedo, A. Barachant, and R. Bhatia, ``Riemannian geometry for
EEG-based brain--computer interfaces; a primer and a review,'' \emph{Brain-Comput.
Interfaces}, vol.~4, no.~3, pp.~155--174, 2017.

\bibitem{lawhern2018} V. J. Lawhern \emph{et al.}, ``EEGNet: A compact convolutional neural
network for EEG-based brain--computer interfaces,'' \emph{J. Neural Eng.}, vol.~15, no.~5,
Art. no.~056013, 2018.

\bibitem{schirrmeister2017} R. T. Schirrmeister \emph{et al.}, ``Deep learning with
convolutional neural networks for EEG decoding and visualization,'' \emph{Hum. Brain Mapp.},
vol.~38, no.~11, pp.~5391--5420, 2017.

\bibitem{roy2019} Y. Roy, H. Banville, I. Albuquerque, A. Gramfort, T. H. Falk, and J.
Faubert, ``Deep learning-based electroencephalography analysis: A systematic review,''
\emph{J. Neural Eng.}, vol.~16, no.~5, Art. no.~051001, 2019.

\bibitem{dempster2021} A. Dempster, D. F. Schmidt, and G. I. Webb, ``MiniRocket: A very fast
(almost) deterministic transform for time series classification,'' in \emph{Proc. ACM SIGKDD},
2021, pp.~248--257.

\bibitem{cortes1995} C. Cortes and V. Vapnik, ``Support-vector networks,'' \emph{Mach.
Learn.}, vol.~20, no.~3, pp.~273--297, 1995.

\bibitem{breiman2001} L. Breiman, ``Random forests,'' \emph{Mach. Learn.}, vol.~45, no.~1,
pp.~5--32, 2001.

\bibitem{chen2016} T. Chen and C. Guestrin, ``XGBoost: A scalable tree boosting system,'' in
\emph{Proc. ACM SIGKDD}, 2016, pp.~785--794.

\bibitem{ke2017} G. Ke, Q. Meng, T. Finley, T. Wang, W. Chen, W. Ma, Q. Ye, and T.-Y. Liu,
``LightGBM: A highly efficient gradient boosting decision tree,'' in \emph{Adv. Neural Inf.
Process. Syst.}, 2017, pp.~3146--3154.

\bibitem{geethapriya2026} M. Geethapriya and K. Pattabiraman, ``Subject-independent
classification of ADHD from pediatric EEG: A leakage-aware deep learning framework with
interpretable neurophysiological mapping,'' \emph{Front. Neurosci.}, vol.~20, 2026, Art.
no.~1925266, doi: 10.3389/fnins.2026.1925266.

\bibitem{maheshwari2025} S. Maheshwari, K. N. V. P. S. Rajesh, V. Kanhangad, U. R. Acharya, and
T. S. Kumar, ``Entropy difference-based EEG channel selection technique for automated detection
of ADHD,'' \emph{PLoS ONE}, vol.~20, no.~4, Art. no.~e0319487, 2025.

\bibitem{bansal2025} J. Bansal, G. Gangwar, M. Aljaidi, A. Alkoradees, and G. Singh,
``EEG-based ADHD classification using autoencoder feature extraction and ResNet with double
augmented attention mechanism,'' \emph{Brain Sci.}, vol.~15, no.~1, Art. no.~95, 2025.

\bibitem{mao2025} Y. Mao, X. Qi, L. He, S. Wang, Z. Wang, and F. Wang, ``Advanced machine
learning techniques reveal multidimensional EEG abnormalities in children with ADHD: A framework
for automatic diagnosis,'' \emph{Front. Psychiatry}, vol.~16, Art. no.~1475936, 2025.

\bibitem{saeb2017} S. Saeb, L. Lonini, A. Jayaraman, D. C. Mohr, and K. P. Kording, ``The need
to approximate the use-case in clinical machine learning,'' \emph{GigaScience}, vol.~6,
no.~5, 2017.

\bibitem{brookshire2024} G. Brookshire \emph{et al.}, ``Data leakage in deep learning studies of
translational EEG,'' \emph{Front. Neurosci.}, vol.~18, Art. no.~1373515, 2024.

\bibitem{varoquaux2018} G. Varoquaux, ``Cross-validation failure: Small sample sizes lead to
large error bars,'' \emph{NeuroImage}, vol.~180, pp.~68--77, 2018.

\bibitem{donoho1995} D. L. Donoho and I. M. Johnstone, ``Adapting to unknown smoothness via
wavelet shrinkage,'' \emph{J. Amer. Statist. Assoc.}, vol.~90, no.~432, pp.~1200--1224, 1995.

\bibitem{efron1993} B. Efron and R. J. Tibshirani, \emph{An Introduction to the Bootstrap}.
New York, NY, USA: Chapman \& Hall, 1993.

\bibitem{clopper1934} C. J. Clopper and E. S. Pearson, ``The use of confidence or fiducial
limits illustrated in the case of the binomial,'' \emph{Biometrika}, vol.~26, no.~4,
pp.~404--413, 1934.

\bibitem{pedregosa2011} F. Pedregosa \emph{et al.}, ``Scikit-learn: Machine learning in
Python,'' \emph{J. Mach. Learn. Res.}, vol.~12, pp.~2825--2830, 2011.

\end{thebibliography}
